\documentclass[12 pt]{article}  % sets the font to 12 pt and says this is an 
                                % article (as opposed to book or other documents)
\usepackage{amsfonts, amssymb}  % math fonts and symbols
\usepackage{amsmath}            % math
\usepackage{enumitem}           % change enumeration bullet styles
\usepackage{graphicx}           % graphics
\usepackage{parskip}
\usepackage{xcolor}             % color
\usepackage{hyperref}           % linkable references
\hypersetup{
    colorlinks=true,
    linkcolor=black,
    filecolor=black,      
    urlcolor=black,
    citecolor=black
    }
\usepackage[nottoc,numbib]{tocbibind}
  
%\usepackage{setspace}          % Together with \doublespacing below allows for
                                % doublespacing of the document

\oddsidemargin=-0.5cm
\setlength{\textwidth}{7in}
\addtolength{\voffset}{-20pt}
\addtolength{\headsep}{25pt}
\renewcommand\thefootnote{[\roman{footnote}]} % Bracketed footnotes

                                        % Macro definitions for the header
\newcommand{\myname}{Danny Nguyen}      % Name
\newcommand{\myid}{OU ID: 113687989}    % Student ID
\newcommand{\mysemester}{Summer 2026}   % Semester
\newcommand{\myclass}{DSA 5900}         % Class
\newcommand{\myassignment}{Mid-Semester Progress Report}  % Assignment designator

% Page headers automatically horizontal fill; Name ID Date Class Assignment Page.No.
\pagestyle{myheadings}
\markright{\myname \hfill \myid \hfill \today \hfill \myclass \hfill \myassignment \hfill }

\newcommand{\eqn}[0]{\begin{array}{rcl}}    % begin an aligned equation - allows
                                            % for aligning = or inequalities.
                                            % Always use with $$ $$
\newcommand{\eqnend}[0]{\end{array} }  	    % end the aligned equation

\newcommand{\qed}[0]{$\square$}     % make an unfilled square the
                                    % default for ending a proof

% Common number sets, must be in math mode
\newcommand{\Naturals}[0]{\mathbb{N}}   % Natural numbers
\newcommand{\Integers}[0]{\mathbb{Z}}   % Integers
\newcommand{\Rationals}[0]{\mathbb{Q}}  % Rationals
\newcommand{\Reals}[0]{\mathbb{R}}      % Reals
\newcommand{\Complex}[0]{\mathbb{C}}    % Complex

\newcommand{\heading}[1]{\noindent\large{\textbf{#1}}}  % for sections
\newcommand{\problem}[1]{\noindent\paragraph{#1:}}      % for problem labels

%\doublespacing         % Together with the package setspace above allows for
                        % doublespacing of the document

% Use these commands for self-assessments
\newcommand{\correct}[1]{\colorbox{green}{\checkmark #1}}   % mark correct
\newcommand{\correction}[1]{{\color{red}\textbf{#1}}}       % write a correction

\title{\myassignment \\ \large \myclass--\mysemester \\ \myname}

\begin{document}

\begin{titlepage}
    \begin{center}
        \vfill
        \Huge \myassignment \\
        \vspace*{0.5in}
        \Large \mysemester -- \myclass -- 1 Credit Hour\\
        \vspace*{0.5in}
        \myname
        \vfill
        \begin{tabular}{rl}
            Faculty Supervisor: & Dr. Matthew Beattie \\
            Thesis Advisor:     & Dr. Charles Nicholson
        \end{tabular}
        \vfill
    \end{center}
\end{titlepage}
\pagenumbering{roman}
\tableofcontents
\listoffigures
\newpage
\pagenumbering{arabic}

\section{Introduction}
Translating natural language into formal predicate logic is a fundamental 
prerequisite for automated reasoning, yet the inherent ambiguity of human 
language makes this process notoriously difficult for current generative 
pipelines. When tasked with outputting structured logical configurations, 
generative LLMs often succumb to hallucination and parsing fragility, rendering 
the resulting logical structures unreliable~\cite{jin2022logicalfallacydetection, 
lalwani2025autoformalizingnaturallanguagefirstorder, lu-etal-2022-parsing,
yang2023couplinglargelanguagemodels}. 

This practicum proposes a shift in methodology: bypassing text-generation in 
favor of parameter-level probing. Rather than forcing models to generate 
structured output, the project embraces the model's existing internal representations. 
The central objective is to explore the model's internal geometry to reveal 
whether logical relationships such as: implication, negation, and 
conjunction; are explicitly encoded within the hidden state vector embeddings 
of Meta's Llama 3.2-3B \cite{alain2018understandingintermediatelayersusing, 
marks2024geometrytruthemergentlinear}. By analyzing these internal activation 
patterns, this practicum aims to determine if verifiable predicate logic can 
be extracted deterministically from the latent space of a compact transformer, 
offering a robust, low-compute alternative to stochastic generative parsing.

% 2. THE FOUNDATION: Transformer Geometry & Latent Representations
\subsection{Transformer Architecture}
The large language model (LLM) under test is Meta's LLaMa 3.2-3B, a 3-billion 
parameter transformer model with 28 transformer head layers, each consisting of 
a grouped-query attention mechanism and Swish feed-forward network
\cite{grattafiori2024llama3herdmodels}. The model is trained on a mixture of 
publicly available and licensed datasets combined with the outputs of the LLaMa 
3.1-8B and 70B models \cite{Meta_2024}.

\subsubsection{Tokenization and Positional Encoding}

The translation of continuous natural language into a structured mathematical 
form begins with tokenization. Tokenization defines a deterministic mapping 
from an input string of characters into a discrete sequence of integer tokens, 
$\mathbf{t} = (t_1, t_2, \dots, t_L)$, drawn from a fixed vocabulary 
$\mathcal{V}$. In modern transformer architectures, tokens do not maintain a 
strict one-to-one correspondence with individual words; rather, subword 
tokenization algorithms (such as Byte-Pair Encoding) decompose rare words into 
sub-word units while mapping common words to single token indices. The 
vocabulary size, $|\mathcal{V}|$ (128K for Llama 3.2-3B), establishes the 
fundamental dimensionality of the discrete input space.

Each discrete token index $t_i \in \mathcal{V}$ is subsequently mapped to a 
continuous vector via an embedding lookup matrix 
$\mathbf{W}_E \in \mathbb{R}^{|\mathcal{V}| \times d_{\text{model}}}$, yielding 
an initial sequence of dense token representations:

\begin{equation}
    \mathbf{x}_{e, i} = \mathbf{W}_E[t_i] \in \mathbb{R}^{d_{\text{model}}}
\end{equation}

where $d_{\text{model}} = 3072$ for Llama 3.2-3B
\cite{grattafiori2024llama3herdmodels,Meta_2024}.

Because standard self-attention operations are inherently 
permutation-invariant—treating inputs as an unordered bag of vectors—the 
sequential order of tokens must be explicitly injected into the model 
representations. Llama 3.2-3B incorporates sequence geometry using Rotary 
Position Embedding (RoPE) \cite{grattafiori2024llama3herdmodels,Meta_2024,
su2023roformerenhancedtransformerrotary}. Unlike additive static positional 
embeddings, RoPE encodes absolute positional information by multiplying key and 
query vectors with a complex-valued rotation matrix in two-dimensional 
sub-chunks:

\begin{equation}
    \mathbf{R}_{\Theta, m}^{d} = \text{diag}
    \left( 
        \mathbf{R}_{\theta_1, m}, 
        \mathbf{R}_{\theta_2, m}, 
        \dots, 
        \mathbf{R}_{\theta_{d/2}, m}
    \right)
\end{equation}

where $\mathbf{R}_{\theta_i, m}$ is a 2D rotation matrix operating on feature 
pairs with angle $m\theta_i$ at token position $m$. This formulation guarantees 
that the inner product between query and key vectors naturally incorporates 
relative distance ($m - n$), preserving syntactic and spatial structure as 
sequence vectors propagate into the primary residual stream.

\subsubsection{Transformer Attention Heads and Feed-Forward Networks}

Within each layer of the decoder stack, the residual state vector 
$\mathbf{h}_l \in \mathbb{R}^{d_{\text{model}}}$ is sequentially processed by 
two core functional components: the Multi-Head Attention block and the 
non-linear Feed-Forward Network (FFN).

\paragraph{Grouped-Query Attention (GQA)}
Standard Multi-Head Attention (MHA) projects hidden states into $N_q$ 
independent Query, Key, and Value heads. During autoregressive decoding, key 
and value tensors across preceding sequence positions are retained in a 
KV-cache to avoid recomputation. While Multi-Query Attention (MQA) reduces 
memory overhead by sharing a single Key and Value head across all Query heads, 
Grouped-Query Attention (GQA) generalizes both approaches by partitioning $N_q$ 
query heads into $N_g$ groups, where each group shares a single Key and Value 
head ($N_{kv} = N_g$) \cite{ainslie2023gqatraininggeneralizedmultiquery}.

Formally, for a group of $G = N_q / N_g$ query vectors 
$\{\mathbf{q}_{g,1}, \dots, \mathbf{q}_{g,G}\}$ assigned to a shared key-value 
pair $(\mathbf{k}_g, \mathbf{v}_g)$, the attention output for head $i$ in group 
$g$ is defined as:

\begin{equation}
    \mathbf{O}_{g,i} = \text{Softmax}
    \left(
        \frac{\mathbf{q}_{g,i} \mathbf{k}_g^T}{\sqrt{d_{\text{head}}}}
    \right)
    \mathbf{v}_g
\end{equation}

Llama 3.2-3B incorporates GQA with $N_q = 24$ query heads mapped across 
$N_{kv} = 8$ key-value heads ($G = 3$), reducing the memory size of the 
KV-cache compared to standard MHA 
\cite{grattafiori2024llama3herdmodels,Meta_2024}.

\paragraph{SwiGLU Feed-Forward Network}
Following attention aggregation, the residual state vector is processed by a 
position-wise Feed-Forward Network (FFN). Rather than using standard 
single-activation MLPs (such as ReLU or GeLU), Llama 3.2-3B implements the 
SwiGLU variant of the Gated Linear Unit (GLU) 
\cite{shazeer2020gluvariantsimprovetransformer,grattafiori2024llama3herdmodels,
Meta_2024}.

SwiGLU applies a gated dual-path linear transformation utilizing the SiLU 
(Swish with $\beta=1$) activation function, where 
$\text{SiLU}(x) = x \cdot \sigma(x)$. Formally, for an input vector 
$\mathbf{x} \in \mathbb{R}^{d_{\text{model}}}$, the transformation is computed 
as:

\begin{equation}
    \text{SwiGLU}(\mathbf{x}) = 
    \left( 
        \text{SiLU}
        \left( 
            \mathbf{x} \mathbf{W}_G 
        \right) \otimes 
        \left( 
            \mathbf{x} \mathbf{W}_U 
        \right) 
    \right) \mathbf{W}_D
\end{equation}

where 
$\mathbf{W}_G, \mathbf{W}_U \in \mathbb{R}^{d_{\text{model}} \times d_{\text{ffn}}}$ 
denote the gate and up-projection weight matrices, 
$\mathbf{W}_D \in \mathbb{R}^{d_{\text{ffn}} \times d_{\text{model}}}$ 
denotes the down-projection matrix, and $\otimes$ represents the Hadamard 
(element-wise) product. In Llama 3.2-3B, the intermediate hidden dimension is 
set to $d_{\text{ffn}} = \frac{8}{3} d_{\text{model}} = 8192$.

\subsubsection{Unembedding, Logit Generation, and Representation Probing Target}

After propagating through the final $L$-th decoder layer, the output vector 
represents the contextualized hidden activation state, 
$\mathbf{h}_L \in \mathbb{R}^{d_{\text{model}}}$, for each token position in 
the sequence.

\paragraph{The Unembedding Transformation and Logits}
To map continuous residual states back into discrete token probabilities for 
autoregressive generation, $\mathbf{h}_L$ is normalized via Root Mean Square 
Normalization (RMSNorm) and projected onto the vocabulary space via the 
unembedding weight matrix 
$\mathbf{W}_{\text{unembed}} \in \mathbb{R}^{|\mathcal{V}| \times d_{\text{model}}}$:

\begin{equation}
    \mathbf{z} = \mathbf{W}_{\text{unembed}} 
    \left( 
        \text{RMSNorm}(\mathbf{h}_L) 
    \right) \in \mathbb{R}^{|\mathcal{V}|}
\end{equation}

The resulting vector $\mathbf{z}$ contains unnormalized log-probabilities, or 
\textit{logits}, corresponding to each token index in the vocabulary 
$\mathcal{V}$. The probability distribution $P(t_i)$ over the vocabulary for 
generating the subsequent token is computed via the Softmax operation:

\begin{equation}
    P(t_i \mid t_{<i}) = \frac{\exp(z_i)}{\sum_{j=1}^{|\mathcal{V}|} \exp(z_j)}
\end{equation}

In standard text generation pipelines, sampling or greedy decoding is performed 
directly on this output probability distribution $P(t_i)$.

% 3. THE PROPOSED SOLUTION: Representation Probing
\subsection{Logical Representation Probing of Intermediate Residual States}
In contrast to standard generative operations that evaluate $\mathbf{z}$ at the 
terminal layer $L$, representation probing targets the intermediate hidden 
activation vectors $\mathbf{h}_l \in \mathbb{R}^{d_{\text{model}}}$ across 
intermediate layers ($l < L$), prior to the unembedding projection. 

Because transformer layers apply additive residual updates 
($\mathbf{h}_l = \mathbf{h}_{l-1} + \text{Attention}(\mathbf{h}_{l-1}) + \text{FFN}(\mathbf{h}_{l-1})$), 
the intermediate residual stream acts as a continuous representation space 
where semantic and structural properties accumulate. Rather than evaluating 
token probabilities from $\mathbf{z}$, linear probes act as diagnostic 
classifiers directly on the intermediate vectors $\mathbf{h}_l$. By training a 
decision boundary parameter vector $\mathbf{w}_p \in \mathbb{R}^{d_{\text{model}}}$ 
on these activations, probing quantifies whether target logical properties 
(such as negation or implication) are linearly encoded within the intermediate 
latent space before the final generation phase.

% 4. METHODOLOGY IN ACTION: Node and Edge Quantifiers

% 5. THE ULTIMATE HORIZON: Automated Logic Graph Extraction
\subsection{Logic Graph Construction}
By validating and assembling these individual classification probes, we establish the core diagnostic machinery required for automated logic graph construction (Figure 1 and Figure 2)... [Explain how atomic probes construct nodes while relational probes draw directed edges, forming a deterministic bridge from natural language to symbolic reasoning].

An example of the training and subsequent inference behavior of a negation
detection probe follows:

\begin{itemize}
    \item \textbf{Input:} Natural Language Text Statement: ``Eli is not soft.''
    \item \textbf{Output:} Log-Odds Score of Negation Presence between the Atomic Elements
\end{itemize}

likewise for an implication detection probe:

\begin{itemize}
    \item \textbf{Input:} Natural Language Text Statement: ``If someone is soft, he is poised and not scared.''
    \item \textbf{Output:} Log-Odds Score of Implication Presence between the Atomic Elements
\end{itemize}

With the same principle mirrored for training classification probes for other 
desired logical relationships such as equivalence, conjunction, and disjunction.

A nominal example of a pipeline utilizing successfully trained probes is 
proposed in Figure~\ref{fig:intro-example-text} with the corresponding 
visualization of the resulting graph in Figure~\ref{fig:intro-example-graph}.

\begin{figure}[h!t]
    \centering
    \includegraphics[width=0.9\textwidth]{figures/intro-example-text.pdf}
    \caption{Pipeline from natural language to logical graph representation.}
    \label{fig:intro-example-text}
\end{figure}

\begin{figure}[h!t]
    \centering
    \includegraphics[width=0.7\textwidth]{figures/intro-example-graph.pdf}
    \caption{Visualization of the logical graph extracted from the natural language example.}
    \label{fig:intro-example-graph}
\end{figure}

\newpage
\section{Objectives}
The overarching goal of this practicum is to develop a robust, parameter-level 
extraction pipeline that bypasses generative LLM bottlenecks to map natural 
language into verifiable logical graphs.
The specific project and individual 
objectives are outlined below:

\subsection{Technical Project Objectives}
\begin{itemize}
    \item \textbf{Latent Space Probing Architecture:} Design and train 
    low-compute, linear classification heads (probes) to extract propositional 
    claims and logical relationships directly from the intermediate residual 
    stream of Meta's Llama 3.2-3B model.
    \item \textbf{Symbolic Pipeline Integration:} Construct a deterministic 
    pipeline that translates extracted latent representations into discrete 
    symbolic tokens without relying on stochastic text-generation or JSON 
    parsing.
    \item \textbf{Graph-Based Fact Verification:} Implement a visualization and 
    verification layer that structures the extracted symbols into a formal 
    logical graph, for human interpretation of the underlying logical 
    structures of natural language.
    \item \textbf{Performance Optimization:} Evaluate the classifier framework 
    to achieve high factual and structural precision while reducing 
    inference-time computational overhead compared to baseline autoregressive 
    generative prompting.
\end{itemize}

\subsection{Individual Learning Objectives}
\begin{itemize}
    \item \textbf{Mechanistic Interpretability Mastery:} Gain deep, practical 
    expertise in analyzing the internal geometry of transformers, specifically 
    concerning layer-wise activation tracking, representation dynamics, and 
    internal knowledge localization.
    \item \textbf{Advanced Representation Engineering:} Master state-of-the-art 
    diagnostic probing and activation patching methodologies using specialized 
    interpretability toolkits.
    \item \textbf{Symbolic Architecture Development:} Develop a rigorous 
    engineering workflow for bridging machine representations (vector 
    embeddings) with classical symbolic paradigms (graph structures).
\end{itemize}

\newpage
\section{Data}

\subsection{Ingestion}
The dataset selected for this project is LogicNLI, an open-source diagnostic 
benchmark hosted on the Hugging Face hub and developed by Tian et al.. LogicNLI 
was specifically created to address growing concerns regarding whether modern 
language models (LMs) truly possess reasoning capabilities or if they merely 
rely on spurious statistical correlations.

As a Natural Language Inference (NLI)-style dataset, LogicNLI is designed to 
isolate First-Order Logical (FOL) reasoning from ``common sense'' and domain 
knowledge. It accomplishes this by utilizing a synthetic composition consisting 
of triplets of ``facts'', ``rules'', and ``statements'' that test a model's 
ability to handle fundamental logical operators; such as conjunction, disjunction, 
negation, implication, equivalence, and universal and existential quantifiers. 
This synthetic approach is highly beneficial for this study because it provides 
a controlled environment to evaluate model performance across four key 
diagnostic perspectives: overall accuracy, robustness to irrelevant information, 
more-hop generalization, and proof-based traceability\cite{tian-etal-2021-diagnosing}.

The data is publicly available and well-documented. For this project the dataset 
is stored locally in a structured directory within the project repository, 
ensuring easy access for preprocessing and analysis.

\subsection{Exploration}
Initial exploratory data analysis revealed a highly uniform, balanced, and 
synthetically controlled dataset structure across its 20,000 entries.
The target \texttt{labels} field contains a near-perfect balanced split across 
its classification categories (Entailment, Contradiction, Self-Contradiction, 
and Neutral) eliminating the risk of majority-class bias during probe 
calibration. A structural analysis of the text fields highlights the specific 
syntactic attributes of the dataset:

\begin{itemize}
    \item \textbf{Facts Array:} This field contains flat, atomic assertions of 
    absolute truth (e.g., \texttt{["Alice is in the garden.", 
    "Bob is hungry."]}).
    These strings exhibit zero logical connectives, 
    maintaining a short average string length of 4 to 6 tokens.
    \item \textbf{Rules Array:} This field holds the propositional conditional 
    formulas (e.g., \texttt{["If Alice is in the garden, then Alice is 
    reading."]} or structures introducing logical equivalence and 
    conjunction (\texttt{"and"})).
    These strings demonstrate higher token 
    complexities and specific syntactic anchors that dictate the relationship 
    edges.
    \item \textbf{Statements and Labels:} The \texttt{statements} field 
    presents the test query, which correlates cleanly to an explicit truth 
    value defined by the formal deduction of the combined rules and facts.
\end{itemize}

An abbreviated example of one raw record from LogicNLI's dataset is shown in 
Figure~\ref{fig:logicnli-example}.

\begin{figure}[h!tbp]
    \centering
    \begin{small}
    \begin{verbatim}
{
  "facts": [
    "Eli is not soft.", "Patricia is civil.", "Broderick is soft.", 
    "Paul is civil.", "Miles is not southern.", "Paul is not scared.", 
    "Ronald is jittery.", "Broderick is not scared.", 
    "Broderick is not poised.", "Paul is not poised.", 
    "Eli is not jittery.", "Eli is not poised."
  ],
  "rules": [
    "If someone is southern, then he is neither jittery nor soft.",
    "If someone is jittery or soft, then he is scared.",
    "As long as someone is soft, he is poised and not scared.",
    "Someone who is not jittery is always both civil and not soft."
    ...
  ],
  "statements": ["Patricia is scared.", "Eli is not soft."],
  "labels": ["entailment", "entailment"]
}
    \end{verbatim}
    \end{small}
    \caption{Example data record from the LogicNLI dataset, illustrating the 
    synthetic structure of facts, rules, and logical statements.}
    \label{fig:logicnli-example}
\end{figure}

Exploratory token length histograms further confirmed tight distribution peaks, 
ensuring that the model's intermediate hidden layers will not suffer from 
sequence padding distortions during batch processing.

\subsection{Preparation}
Because the linear probes require distinct feature-target pairs to map out 
Llama's internal vector space, substantial preprocessing and feature 
engineering were executed to transform the raw dataset into a token-level 
sequence classification task.
No missing data or null fields were detected due 
to the synthetic nature of LogicNLI, bypassing the need for standard 
statistical imputation.
The data preparation workflow consists of two critical pipeline stages:
\begin{enumerate}
    \item \textbf{Rule-Based Target Transformation and Fact-Rule Matching:} A 
    custom semantic parsing script was built to map the structural elements of 
    the environment. Within the \texttt{rules} array, regular expressions and 
    token-dependency parsing isolate the relational connectives, explicitly 
    identifying whether a rule represents an implication ($\rightarrow$) or a 
    conjunction ($\land$). Concurrently, a lookup matrix maps these rules 
    directly back to the atomic strings inside the \texttt{facts} array. This 
    enables the framework to identify exactly which baseline facts are being 
    operated on by any given rule, providing clear structural targets for the 
    relational edge probes.
    \item \textbf{Contradiction-Driven Negation Curation:} Because explicit 
    negation ($\neg$) is frequently represented implicitly across text pairs 
    rather than through simple keyword matching, a target-generation strategy 
    was implemented leveraging the dataset's logical symmetries. The pipeline 
    filters entries where the target \texttt{label} is explicitly marked as a 
    \texttt{"contradiction"}. By aligning the queried \texttt{statement} 
    against its corresponding asset in the \texttt{facts} array under this 
    specific label, the pipeline mathematically isolates the exact token span 
    undergoing negation. This contradiction-driven mapping provides a 
    high-precision, curated subset of explicit positive-to-negative vector 
    transformations.
    \item \textbf{Token Span Realignment:} Because character indices shift once 
    strings are converted into transformer sub-tokens, the Hugging Face 
    tokenizer was configured with the \\\texttt{return\_offsets\_mapping=True} 
    parameter. This maps the character spans generated during the rule-parsing 
    and contradiction-alignment steps directly onto the sequence tensor 
    positions of Meta's Llama 3.2-3B model. This preparation ensures that when 
    forward hooks intercept the hidden states of layers 8 through 16, the 
    mean-pooling functions grab the precise vector slices corresponding to the 
    isolated logical components.
\end{enumerate}

\newpage
\section{Methodology}

\subsection{Techniques}
To extract and verify the logical structures embedded within Meta's Llama 3.2-3B 
model, this practicum utilizes a mechanistic interpretability framework centered 
on diagnostic linear ``probing''.

\paragraph{Transformers}
In the standard Transformer architecture, input sequences are processed through 
a stack of $L$ sequential transformer blocks. For any given token at position 
$i$ and layer $l$, the model generates a hidden state activation vector 
$h_i^{(l)} \in \mathbb{R}^d$. This vector serves as the distributed, 
high-dimensional representation of the token's contextual state, encoding the 
compositional hierarchy of syntactic dependencies and semantic features derived 
from the preceding attention and feed-forward operations.

These activations persist within the ``residual stream'', an architectural design 
that facilitates the additive accumulation of information across the network's 
depth. In this framework, each transformer block $l$ performs an incremental 
update $h_i^{(l+1)} = h_i^{(l)} + \text{sublayer}(h_i^{(l)})$, ensuring that 
the residual stream serves as a continuous conduit for information propagation. 
This project posits that this latent space encodes propositional logical 
structures—such as implications, negations, and conjunctive relations—as 
geometric properties within the manifold of the hidden state embeddings, 
independent of, and prior to, the final logit projection onto the vocabulary 
space.

By applying diagnostic linear classification ``probes'' to these frozen activation 
vectors, the linear separability of these logical relations can be formally 
tested. Specifically, if a discrete logical relation $Y \in \{0, 1\}$ is 
decodable via a weight matrix $W$, it provides empirical evidence that the 
information is explicitly formatted in a linearly recoverable configuration 
within the model's internal representation of the input.

\paragraph{Diagnostic Linear Classification Probing}
To facilitate the extraction of latent logical structures, this project 
implements diagnostic linear classification, utilizing the methodology 
established by Alain and Bengio (2018). The core principle of linear probing 
involves treating the internal activation vectors of a frozen transformer model 
as a fixed feature space. By training a simple, low-capacity linear model—a 
``probe"—to map these high-dimensional hidden state vectors to discrete logical 
class labels, the linear separability of internal concept representations within 
the model's latent geometry can be empirically assessed.

In contrast to end-to-end fine-tuning, which modifies the parameters of the base 
architecture, linear probing relies on the gradient descent optimization of a 
shallow classifier using frozen activation outputs. High classification accuracy 
on held-out test sets provides empirical evidence that the target logical 
relations are explicitly formatted in a linearly recoverable configuration at 
the probed layer. Conversely, insufficient classification performance suggests 
that the information is either absent from the latent representation or encoded 
within a non-linear manifold that necessitates a more complex, high-capacity 
decodable mapping.

\subsubsection{Algorithmic Formulation and Modeling}
To quantify the presence of logical relations within the transformer's latent 
space, this project employs shallow linear classification probes. A ``hidden 
state activation vector" is the high-dimensional numerical output of a 
transformer block at a specific layer, representing the model's contextual 
understanding of an input token. A ``shallow linear probe" is a low-capacity 
classifier (specifically, logistic regression) designed to identify patterns in 
these vectors without modifying the underlying transformer parameters. The 
``discrete logical class $Y$" refers to the target category—such as the presence 
of an implication or a negation—that the probe is trained to predict from the 
hidden state activation vector.

The classification process is optimized via Logistic Regression with $L_2$ 
regularization to prevent over-fitting to the training samples. Given a dataset 
of hidden state activation vectors $X \in \Reals^{d}$ extracted from a specific 
intermediate transformer layer $l$ (where $d = 2048$ for Llama 3.2-3B), the 
probe learns a weight matrix $W$ and bias vector $b$ to map the input hidden 
state $X$ to the discrete logical class $Y$:

$$\hat{Y} = \sigma(W \cdot X + b)$$

where $\sigma$ denotes the sigmoid function. The weight matrix $W$ acts as a 
learned linear decision boundary within the $d$-dimensional embedding space. 
By optimizing $W$ via gradient descent, the probe identifies the hyperplane 
that optimally separates the hidden state vectors corresponding to the presence 
versus the absence of a logical relation. The output $\hat{Y}$ represents the 
predicted probability of the logical class. This mapping is derived from the 
linear log-odds (logit), where $W \cdot X + b$ defines the decision boundary in 
the latent space.

Once trained, $W$ serves as a diagnostic tool for logical decoding. During 
evaluation, new hidden state vectors are projected through the learned $W$ and 
$b$; the resulting scalar output $\hat{Y}$ indicates the model's internal 
assignment of a logical class to that specific token position. High 
classification accuracy on unseen test data demonstrates that the logical 
relation $Y$ is not merely implicit but is explicitly organized as a linearly 
separable feature in the model's internal representation, thereby validating the 
model's capacity to internalize structured logical operations. Conversely, poor 
classification performance suggests the absence of a linearly recoverable 
representation of the target logical relation.

\subsubsection{Train-Test Splits and Validation Strategy}
The experimental validation utilizes a predefined train/dev/test split strategy, 
leveraging the existing partitions within the LogicNLI benchmark to ensure a 
rigorous evaluation of the latent logical representations. The dataset comprises 
a total of 20,000 samples distributed as 8,000 for training, 1,000 for 
development (dev), and 1,000 for testing. To minimize bias, class distributions 
across all splits are balanced, ensuring that the probes are trained on an equal 
representation of "entailment" and "contradiction" (or the presence versus 
absence of specific logical relations).

Individual probes are instantiated for each target logical category to isolate 
their distinct representation within the latent space. Separate models are 
trained for the following primitives:

\begin{itemize}
    \item \textbf{Propositional Connectives:} Implication ($\rightarrow$), Negation ($\neg$), and Conjunction ($\land$).
    \item \textbf{Atomic Terms:} Predicate and entity extraction (e.g., \textit{Studies(x)}, \textit{Passes(x)}).
\end{itemize}

The validation strategy adheres to the following pipeline:

\begin{itemize}
    \item \textbf{Calibration Split:} Probes are trained using activation 
    vectors extracted strictly from the 8,000 records of the \texttt{train} split. 
    Baseline performance is established by comparing the linear probe results 
    against a majority-class classifier, ensuring that any successful extraction 
    is indicative of meaningful latent organization rather than heuristic 
    guessing.
    \item \textbf{Model Selection and Layer Auditing (Dev Split):} The 
    validation framework utilizes the \texttt{dev} split to audit the model 
    layer-by-layer. Probes are independently trained and evaluated across layers 
    8 through 16 to map the trajectory of logical synthesis. This specific layer 
    range is selected based on empirical precedents in transformer 
    interpretability, which suggest that mid-to-late layers frequently encode 
    the most sophisticated contextual and relational information prior to the 
    final output layers~\cite{marks2024geometrytruthemergentlinear}. The layer 
    yielding the highest F1-score on the dev set is designated as the optimal 
    extraction layer for that specific primitive.
    \item \textbf{Generalization Verification (Test Split):} Final model 
    performance is reported strictly using vectors from the unseen \texttt{test} 
    split. Out-of-sample generalization is quantified using Precision, Recall, 
    and the $F_1$-score relative to the ground-truth target structures, 
    providing a robust measurement of whether the model has truly internalized 
    the logical relationship as a generalizable feature.
\end{itemize}

\subsection{Process Validation}
The choice to implement a parameter level probing methodology rather than 
standard auto-regressive token generation was refined through consultations with 
the thesis advisor about the application of modern data science techniques to 
automated fact checking. These discussions led to delving into the space of 
mechanistic interpretability research largely around the problems of these token 
generation models when applied to logical reasoning tasks. While auto-regressive 
prediction is the standard paradigm for language modeling, it is often 
suboptimal for diagnostic tasks. Auto-regressive loops generate output 
sequentially, compounding potential errors through stochastic sampling; when an 
LLM is forced to output structured logical text, it frequently suffers from 
hallucination where the model prioritizes syntactic fluency over logical 
consistency~\cite{cheng2025empoweringllmslogicalreasoning}.

Relying on generative output to infer reasoning processes is problematic because 
it conflates the model's ability to reason with its ability to articulate that 
reasoning in natural language~\cite{cheng2025empoweringllmslogicalreasoning}. 
By shifting the focus to internal parameter-level mechanics, the research path 
was restructured to prioritize the extraction of latent representations before 
they are subjected to the non-linear transformations of the final decoding heads. 
This approach adheres to the principle that internal hidden states represent a 
more stable, pre-output manifold of the model's learned logic.

Regarding the project objectives, the focus is placed on enhancing 
interpretability and maintaining computational feasibility rather than claiming 
a final solution to the ``black-box problem''. The primary aim is to investigate 
whether specific logical labels are decodable from internal representations, 
which may offer a more nuanced view of the model's latent geometry. By mapping 
these extracted geometries directly to nodes and edges in a symbolic graph, 
this project tests the hypothesis that propositional logic is organized as 
linearly separable features within the model's high-dimensional space.

This methodology also aligns with the research requirement for computational 
efficiency. By utilizing a compact 3B-parameter model and shallow classification 
heads, the project bypasses the need for massive, closed-source APIs or 
computationally expensive fine-tuning loops. This engineering path ensures that 
the extraction and verification architecture maintains high structural precision 
while operating within the memory and throughput constraints of consumer-grade 
hardware (such as an NVIDIA RTX 4070 Ti). This validation process establishes a 
reproducible, low-overhead framework capable of bridging deep learning vector 
spaces with verifiable symbolic graph structures.

\newpage
\section{Results}
% This section is optional for this progress report.  It corresponds to the 
% Evaluation phase of the data science method.
TBD. I hope I have enough time try single-layer and multi-layer probes, and a 
couple different open models to see how results change by model size.

\newpage
\section{Deliverables}
% This section is optional for this progress report.  It corresponds to the 
% Deployment phase of the data science method.
TBD

\cleardoublepage
\bibliographystyle{plain}
\bibliography{../refs.bib}
\end{document}